\addtocontents{toc}{\protect\Needspace{3\baselineskip}}
\originalchapter{AdS 几何与对偶字典}\label{ch:AdS}
\originalsection{径向方向与边界}
AdS$_{d+1}$ 是双时间嵌入空间中的双曲面 $-X_{-1}^2-X_0^2+\sum_{i=1}^dX_i^2=-R^2$. 全局度规可写
\begin{equation}
\dd s^2=R^2[-(1+r^2)\dd t^2+\dd r^2/(1+r^2)+r^2\dd\Omega_{d-1}^2].
\end{equation}
通常使用覆盖空间消去全局时间周期所产生的闭时线. Poincaré 坐标中 $\dd s^2=R^2(\dd z^2+\eta_{\mu\nu}\dd x^\mu\dd x^\nu)/z^2$, $z>0$, 只覆盖 Lorentzian AdS 的一个区域. 边界 $z\to0$ 的整体缩放 $z,x\to\lambda(z,x)$ 是体等距变换, 对边界正是尺度变换. 因而体中较大 $z$ 对应边界较长距离, 是 UV/IR 联系的几何入口.

AdS 的边界度规仅定义到共形因子. 固定 $z=\epsilon$ 作截止时, 坐标格点尺度 $\epsilon$ 与边界紫外尺度对应. 全局边界是 $S^{d-1}\times\mathbb R$, 球半径提供有限尺度, 与平直边界无限平面不同. 这种联系并不意味着局域体点严格等于一个边界局域算符.

\originalsection{D3 膜与参数映射}
$N$ 个 D3 膜的远区开放弦低能理论是四维 $\mathcal N=4$ Yang--Mills; 相应超引力解在近视界区域出现 AdS$_5\times S^5$. 两种低能描述提示对偶猜想. 常用约定为
$R^4=4\pi g_sN\alpha'^2$, $g_{\rm YM}^2=4\pi g_s$, 因而 $R^4/\alpha'^2=\lambda$.
不同规范场迹规范化可让 $g_{\rm YM}$ 的系数变化, 使用公式必须成套保持一致.

经典引力近似需要 $N\gg1$ 抑制量子环, $\lambda\gg1$ 使 $R\gg\sqrt{\alpha'}$ 抑制弦尺度修正, 且相应 $g_s$ 小. 它通常对应边界强耦合, 并非边界普通微扰展开. 对偶的一般成立、有限 $N$ 与有限耦合的所有修正, 不由这两个近似条件证明.

\originalsection{生成泛函}
对每个边界算符 $\mathcal O$, 指定相应体场 $\phi$, 非正规化边界系数 $\phi_0$ 作为源. 欧氏字典写为
\begin{equation}\label{eq:GKPW}
\left\langle\exp\left(\int\phi_0\mathcal O\right)\right\rangle_{\rm CFT}
=Z_{
\rm bulk}[\phi_0]\simeq\exp[-S_{E,\rm ren}^{\rm on-shell}[\phi_0]].
\end{equation}
右侧须指定内部正则性、边界条件与局域反项, 否则作用量通常发散. 取源泛函导数得到关联函数; 两点函数来自自由二次核, 高点函数来自相互作用的 Witten 图. 这是使用对偶字典进行计算的规则, 不是把左侧路径积分从零证明等于右侧.
